\thispagestyle{plain}
\begingroup\centering
{\fontsize{9.5}{10.7}\selectfont Manuscrito de recopilación para transferencia académica\par}
\vspace{3mm}
{\fontsize{17}{20.5}\selectfont\bfseries\Titulo\par}
\vspace{6pt}
{\fontsize{11.5}{14}\selectfont\Autores\par}
\vspace{3pt}
{\fontsize{8}{10}\selectfont Coautoría sin prelación de contribución\par}
\vspace{3pt}
{\fontsize{8}{10}\selectfont Integración y recopilación documental generadas por ChatGPT - Astra.\par}
\par\endgroup
\vspace{5pt}
{\centering\bfseries Resumen\par}
\vspace{2pt}
\begingroup
\fontsize{9.5}{10.7}\selectfont
\setlength{\parskip}{1.5pt plus .5pt minus .5pt}
\noindent
Se generan coinductivamente \(\pi,\varphi,e\) de forma conjunta y a profundidad arbitraria, se
deriva la constante de estructura fina \(\alpha\) y se construyen la
coordenada electrónica, su representación de espín \(1/2\) y sus
componentes de helicidad. La construcción relaciona la determinación de
esos valores con la conservación de las estructuras que los producen.
Su fundamento es la \emph{Holografía Modular Triádica} (HMT):
\emph{All Possible Paths} (APP) es el soporte toroidal discreto
\((\mathbb Z/9\mathbb Z)^2\), equipado con el grafo de Cayley aditivo
de generadores \(\pm(1,0),\pm(0,1)\) y con las evaluaciones suma y
producto; TRIT es la firma ternaria orientada \(\{-1,0,+1\}\), que
distingue tres regímenes cuadráticos; el \emph{Topological Phase Kernel}
(TPK) compone selección de fase, transporte y actualización de memoria.

Las 104\,976 condiciones iniciales orientadas producen 468 emisiones y 243
secuencias ternarias visibles. En este dominio se determinan regiones de
clausura, propagación y autoescala, cuyas realizaciones arquimedianas son
\(\pi,e,\varphi\). El refinamiento por cilindros compatibles permite obtener
cualquier prefijo finito mediante un cálculo finito. La holonomía nonádica
conserva la fase mientras incrementa la memoria; su representación aperiódica
explicita la coexistencia de periodicidad modular y prolongación ilimitada.

El registro integral dodecafásico \(K\) y su realización racional periódica
\(\kappa\) admiten reconstrucciones reversibles. Su composición con los registros
regionales determina \(\alpha\) por normalización posicional. La vía analítica
correlacionada la representa como raíz simple única de una prolongación
del resolvente de vacancias, con recurrencia y clase geométrica explícitas.
La compatibilidad a toda profundidad demuestra la igualdad de ambas
representaciones del mismo estado. La evaluación escalar y la resolución de incidencia
se articulan mediante los registros orientados que ambas conservan.

La construcción electrónica reúne el centro de APP, la vacancia orientada,
una representación matricial de los tres regímenes y la transferencia entre
las evaluaciones areal y volumétrica. Del bloque electrónico se construyen
los generadores de espín \(1/2\) y los proyectores direccionales; su interpretación
como helicidad declara la dirección de propagación y conserva la evaluación
escalar, la orientación de carga y la memoria de la ruta como datos distintos.
Dos lectores del retorno determinan
los coeficientes de su composición completa. La trayectoria central visita
exactamente \(27\) celdas; sus historias orientadas se distinguen mediante
la coordenatización integral reversible. La descomposición armónica del retorno
produce los modos \(2,6,10\), cuyas potencias \(64,16,64\) recuperan
los coeficientes \(80\) y \(6\) del corrector electrónico.
La escala de acción y la
realización dimensional se exponen antes del contraste con CODATA 2018 y
2022. La segunda realización de la estructura conserva la incidencia de
frontera: diseños de Witt, códigos de Golay, retículo de Leech y álgebra
de vértices \(V^\natural\), con su grupo de automorfismos, el Monstruo.
Las realizaciones numérica y excepcional proceden de la misma estructura;
sus operaciones locales y sus límites compatibles conservan los registros
que enlazan ecuaciones, valores e incidencia.
El registro \(K\) selecciona además la dirección de la reducción
\(11\to10\). Los orbifolds de órdenes \(2\) y \(3\) proporcionan
presentaciones isomorfas de \(V^\natural\), con transporte de sus
automorfismos y trazas graduadas. El isomorfismo de las pantallas
cocientadas entrelaza la dualidad de radios recíprocos, conservando
los datos de reducción de \(K\) y la pantalla reticular.

% BEGIN S27_FRAMING
La comparación de las estrellas de \(B_K\), \(B_\varphi\) y
\(B_{\rm APP}\) distingue sus doce completados mediante una polarización
común y sitúa sus firmas en la clasificación de las \(495\) cuaternas.
El marco ordenado \((H^-_\alpha,H_e,H_\varphi,H_{\rm APP})\) tiene
estabilizador trivial en el grupo incidencial construido. El cociclo
de acarreo explicita, a su vez, la información de los cocientes enteros
que complementa los residuos de las columnas regionales. Sus leyes de
composición y transporte conservan el levantamiento aritmético asociado
a la lectura numérica, mientras las estrellas y sus estabilizadores
determinan la organización de la lectura de frontera.

% END S27_FRAMING
\vspace{3pt}
La monodromía del transporte y el modelo matemático de cristal temporal
aperiódico articulan retorno modular, avance de memoria y los tres
regímenes de la exponencial. La década determinantal y el carácter de
acción determinan la sección correspondiente a la constante de Planck
reducida, \(\hbar\), que interviene en la composición electrónica.



% BEGIN R27_FRAMING_SUMMARY_ES
La precisión de esta cadena se desarrolla desde la aritmética del emisor:
se determina su alfabeto decimal de \(42\) bloques, se exhiben los empalmes
entre emisiones regionales y se separan memoria ordenada y distribución
marginal. La selección de frontera distingue saturación, neutralidad dual
y orientación. Las leyes de refinamiento APP, la subretícula de índice
dos, la suma nonádica con acarreo y la correlación cuadrática explicitan
las operaciones intermedias. En la realización de nivel tres se obtiene,
mediante una identidad de series formales, una divisibilidad uniforme
óptima por \(3^9\).
% END R27_FRAMING_SUMMARY_ES

\noindent\textbf{Palabras clave:} derivación estructural de constantes fundamentales;
Holografía Modular Triádica; generación coinductiva; constante de estructura fina;
estructura algebraica del electrón; holonomía; incidencia de frontera;
constante de Planck reducida; monodromía; cristal temporal aperiódico;
simetrías excepcionales; Moonshine.
\par\endgroup
